Tomographic measurements of gravitational lensing with different lens and source redshift distributions contain crucial information about the relative expansion rate of the Universe, and hence dark energy. 
%
% MSD
While this technique is well-established in weak lensing, its application to the strong lensing regime has traditionally focused on rare Double Source Plane Lenses (DSPLs). However, DSPLs are fundamentally limited by the Mass-Sheet Degeneracy (MSD)—a systematic uncertainty underexplored in previous forecasting literature—as well as their intrinsic scarcity.
%
%
% PMSPLs
To overcome these challenges, we introduce \emph{Pseudo} Double-Source Plane Lenses 
(PDSPLs): pairs of independent single-source plane lenses with self-similar deflectors, thus generalizing the DSPL formalism to the $\sim\mathcal{O}(10^5)$ systems expected from LSST, Euclid, and Roman. Unlike true DSPLs, PDSPLs are free from the intermediate source mass problem by construction, eliminating the associated secondary MSD and the need for multi-plane ray tracing. We incorporate the MSD into a hierarchical forecasting framework and demonstrate that this degeneracy severely degrades constraints from small DSPL samples, motivating the statistical approach of PDSPLs.

%
% Results
We forecast constraints on the dark energy equation of state under a Flat 
$w_0w_a$CDM cosmology. The LSST 10-year photometric sample alone achieves 
$\sigma(w_0) \sim 0.43$, while simultaneously constraining the mass-sheet parameter 
and deflector power-law slope to $\sim2\%$. Adding a prior $\mathcal{N}(0.3, 0.05)$ 
on $\Omega_{\rm m}$—simulating combination with external probes such as CMB, BAO, or 
Type Ia supernovae—tightens this to $\sigma(w_0) \sim 0.27$, competitive with current 
Stage III weak lensing analyses. Notably, the large photometric LSST sample 
outperforms the smaller spectroscopic 4MOST subsets, confirming that statistical 
volume dominates over per-pair precision.